\begin{pastexam}[label={px:gw3-103-ds-04}]{關務三等資訊處理 103 年 第四題}
（申論；20 分）

請使用 C 或 Java 語言寫一副程式 void FindMinMax(int [] A, int Min, int Max)，此副程式對一個長度為 10 的整數陣列
A[0:9]，最多花費 15 次的數值比較運算，尋找陣列中的最小值及最大值，並分別存入 Min 及 Max。（20 分）\\
（注意：請加註解說明程式碼作法）
\end{pastexam}

\begin{pastsoln}
答案：兩兩一組比較：先比 $A[0]$、$A[1]$ 定出初始的最小、最大值（$1$ 次），之後每一對 $A[2k]$、$A[2k+1]$ 先互比，小的
和目前最小值比、大的和目前最大值比（每對 $3$ 次），共 $1 + 4 \times 3 = 13$ 次 $\le 15$。

\textbf{題目的參數。} C 和 Java 的 \texttt{int} 參數都是傳值，副程式裡改 \texttt{Min}、\texttt{Max} 不會影響呼叫端，照題目
的宣告寫是拿不回結果的。 所以 C 版本改用指標 \texttt{int *Min, int *Max}（Java 可以改成回傳一個長度 $2$ 的陣列，或把結果
存在物件的欄位），做法不變。 下面「比較」指陣列元素之間的數值比較，迴圈變數的比較不計（題目的意思是資料的比較次數）。

\textbf{想法}（資結書第~16.1~節）。 分別找最小值和最大值要 $9 + 9 = 18$ 次，超過 $15$。 改成一次拿兩個數：兩個數先互比
（$1$ 次），較小的那個才有可能是最小值、較大的才有可能是最大值，所以小的只和目前的最小值比、大的只和目前的最大值比
（$2$ 次）。 每兩個數花 $3$ 次而不是 $4$ 次。
\begin{minted}[fontsize=\small]{c}
/* 在 A[0..9] 中找最小值和最大值，結果存入 *Min、*Max。
   元素之間的比較共 1 + 4*3 = 13 次。 */
void FindMinMax(int A[], int *Min, int *Max) {
    int i, small, big;
    /* 第 1 次比較：前兩個數決定初始的最小值、最大值 */
    if (A[0] < A[1]) { *Min = A[0]; *Max = A[1]; }
    else             { *Min = A[1]; *Max = A[0]; }
    /* 其餘 8 個數兩兩一組：(A[2],A[3]), (A[4],A[5]), (A[6],A[7]), (A[8],A[9]) */
    for (i = 2; i < 10; i += 2) {
        if (A[i] < A[i + 1]) {           /* 第 1 次：同一組的兩個數互比 */
            small = A[i];  big = A[i + 1];
        } else {
            small = A[i + 1];  big = A[i];
        }
        if (small < *Min) *Min = small;  /* 第 2 次：小的只可能更新最小值 */
        if (big > *Max)   *Max = big;    /* 第 3 次：大的只可能更新最大值 */
    }
}
\end{minted}
\textbf{比較次數}：第一組 $1$ 次；迴圈跑 $i = 2, 4, 6, 8$ 共 $4$ 輪，每輪恰好 $3$ 次。 合計 $1 + 12 = 13$ 次，不論資料
如何都是 $13$ 次，符合「最多 $15$ 次」。 一般 $n$ 個數（$n$ 為偶數）這樣做是 $1 + 3(n-2)/2 = \tfrac{3n}{2} - 2$ 次，
和資結書定理~16.1 的 $\lceil 3n/2 \rceil - 2$ 相同；資結書第~16.1~節也用對手論證說明這已經是最少，$n = 10$ 時正好是 $13$。

\textbf{正確性}：迴圈不變式「處理完 $A[0..i-1]$ 後，\texttt{*Min}、\texttt{*Max} 是這一段的最小值、最大值」。 一開始
$i = 2$ 成立。 加入一組 $A[i], A[i+1]$ 時，新的最小值是 $\min(\texttt{*Min}, A[i], A[i+1]) = \min(\texttt{*Min}, \text{small})$，
因為 big $\ge$ small 不會是最小值；最大值同理。 所以迴圈結束（$i = 10$）時得到整個陣列的最小、最大值。

呼叫方式：\texttt{int mn, mx; FindMinMax(A, \&mn, \&mx);}。 Java 寫法相同，例如 \texttt{static int[] FindMinMax(int[] A)}
回傳 \texttt{new int[]\{min, max\}}。
\end{pastsoln}
